% ------------------------------------------------------------
% Page 70
% ------------------------------------------------------------

\begin{center}
    {\Large\bfseries LE MOULIN D’AYRES}
\end{center}

\vspace{0.55cm}

Nous n’avons que de rares documents nous permettant de retracer les origines du moulin :

\vspace{0.35cm}

A-t-il été construit vers les 11° ou 12° siècle comme beaucoup de moulins en Languedoc ?
Nous l’ignorons. Par contre la photocopie d’un acte retrouvé aux Archives de Mende par
Mme Jeanjean-Sanch, datant de 1574, sous Henri III, vient de m’être remis. Il s’agit d’un
contrat de location du moulin d’Ayres par le seigneur ou le sieur de Fontanille à un certain
Voller ou Valles en cours de déchiffrage en Septembre 2008.

\vspace{0.35cm}

Nous savons également qu’en 1736 il était en pleine activité comme moulin à farine, et il est
répertorié, et lui seul, dans les vallées avec celui de Malbois dans le bourg, sur la carte de
Cassini (1750) mais en 1.900 il n’y avait plus, au Moulin d’Ayres, ni roue, ni meules, !
ni aucun vestige à l’exception :

\vspace{0.3cm}

1/ Du canal d’amenée d’eau encore visible ; il débouche dans la cave de l’aile Sud des
bâtiments actuels ; canal bâti en pierres taillées et ajustées.

\vspace{0.3cm}

2/ D’un acte de « vente » passé devant P.~Berthezène, notaire à Meyrueis à un certain Jean
Rozier, meunier.

\vspace{0.35cm}

A l’époque le moulin ne comportait vraisemblablement qu’un petit bâtiment : le moulin
proprement dit : les meules au niveau de la Jonte et au dessus 2 pièces pour habitation (?).
Il se pourrait que le Rouilladou avec son foulon ait déjà existé.

\vspace{0.35cm}

Le Moulin d’Ayres a donc appartenu à \textbf{Jean Rozier} et à ses descendants jusqu’en 1820.
Le dernier \textbf{Gabriel Rouzier} a vendu le domaine et le moulin toujours en activité, à \textbf{Jean
Louis Laget} comme en fait foi l’acte suivant :

\vspace{0.25cm}

\begin{flushleft}
\hspace{1cm}
\parbox{0.82\textwidth}{
\textit{
« Le 20 novembre 1820 le sieur \textbf{Gabriel Rouzier}, meunier.... a procédé à la vente pure,
simple, à jamais valable, irrévocable avec promesse de faire jouir en paix et d’être en toute garantie
de droit ou d’érection en faveur du sieur \textbf{Jean-Louis Laget}, pour le prix de 4.900 F, un
petit doncs de domaine...
Indépendamment du prix ci-dessus stipulé le sieur Laget demeure chargé de payer à M. de
Montauran, ou à ses successeurs la rente ou pension foncière annuelle perpétuelle de 45,60
Francs ...
Le domaine dispose de 2 jardins indépendants : l’un au devant de la maison d’habitation et se
tenant au chemin qui conduit à la dite maison, l’autre derrière la dite maison et moulin se tenant
d’un côté avec l’aire de l’autre à la dite rivière de Jaonte et du chemin... »}
}
\end{flushleft}

\vspace{0.25cm}

\noindent
\textbf{Jean Louis Laget} était propriétaire foncier à Meyrueis ; a-t-il affermé le moulin sans être
lui même le maître-meunier ? C’est possible. En effet nous savons qu’en 1844 le moulin
était toujours en activité et le meunier, fermier également, à Ayres s’appelait \textbf{Trinchard}.
C’est sans doute vers les années 1850 que les bâtiments furent agrandis pour devenir tels,
ou à peu près, que nous les connaissons aujourd’hui. Ce serait donc \textbf{Jean Louis Laget}
qui aurait financé ces travaux importants.
